% Generated draft: requires mathematical and source review.
\documentclass[11pt]{article}
\usepackage[margin=25mm]{geometry}
\usepackage{fontspec}
\setmainfont{Latin Modern Roman}
\usepackage{amsmath,amssymb,mathtools,unicode-math}
\setmathfont{Latin Modern Math}
\usepackage{enumitem,needspace,xurl,hyperref}
\setlength{\parindent}{0pt}
\setlength{\parskip}{5pt}
\setlength{\emergencystretch}{3em}
\newcommand{\N}{\mathbb{N}}
\newcommand{\Z}{\mathbb{Z}}
\newcommand{\Q}{\mathbb{Q}}
\newcommand{\R}{\mathbb{R}}
\newcommand{\C}{\mathbb{C}}
\newcommand{\F}{\mathbb{F}}
\newcommand{\A}{\mathbb{A}}
\newcommand{\PP}{\mathbb{P}}
\newcommand{\E}{\mathbb{E}}
\newcommand{\eps}{\varepsilon}
\newcommand{\dd}{\,\mathrm{d}}
\begin{document}
\section*{UnsolvedMath ID: 11000122}
\textbf{Source problem code:} catalogue entry
\subsection{Definitions and mathematical statement}
Let $N$ be the nonorientable surface intended by the source, as identified by the registry summary, and put $G=\operatorname{Mod}(N)$, the group of homeomorphisms modulo the relevant isotopies. The excerpt does not give the genus, boundary or puncture conventions for $N$.

Problem 3.4 first asks whether every homomorphism $\varphi:G\to G$ with $[G:\varphi(G)]<\infty$ is an automorphism. In symbols, is
\[
 \forall\varphi\in\operatorname{Hom}(G,G),\qquad
 [G:\varphi(G)]<\infty\Longrightarrow
 \varphi\text{ is bijective}?
\]
Injectivity is not an assumption, and finite-index image need not mean surjectivity.

The source then asks for the analogue when the domain is a finite-index subgroup $H\leq G$, with $\varphi:H\to G$ still having finite-index image. What ``automorphism'' should mean in this second formulation is not specified: an isomorphism onto its image, an isomorphism $H\to G$, or the restriction of an ambient automorphism are distinct conclusions. In particular, the inclusion of a proper finite-index subgroup has finite-index image but is not surjective onto $G$, so the literal surjectivity version cannot be intended without further restrictions.

The registry suggests using nonorientable curve complexes and expects homeomorphism-induced rigidity, but gives no exact result establishing either formulation. That expectation does not resolve the missing surface hypotheses or the finite-index-domain conclusion.
\subsection{Short English statement}
Must a finite-index-image endomorphism of a nonorientable surface mapping class group be an automorphism? Clarify the corresponding rigidity statement when its domain is only a finite-index subgroup.
\subsection{Sources}
\begin{itemize}
\item Local registry: \url{data/unsolvedmath/problems.json}, permanent UnsolvedMath ID $11000122$, exact source code \texttt{catalogue entry}.
\item Farb, \emph{Problems on Mapping Class Groups and Related Topics}, Problem 3.4, PDF page 98: \url{https://www.math.uchicago.edu/~farb/papers/mcgbook.pdf}. Supplied access date: 2026-07-29; extraction: pdf-font-structure; repeated local numbering disambiguated by page among 329 headings.
\item The registry mentions nonorientable curve-complex methods and an affirmative expectation without a precise primary citation or individual attribution.
\end{itemize}
\end{document}
